\chapter{Building Production-Grade Projects}
\label{ch:projects}

This is the proof-of-work layer --- the $S$ (Social Proof) term in the outreach equation from
Chapter~\ref{ch:outreach-engineering}. Everything in this Part draws on two aligned sources: the project-related
section of the remote-jobs playbook, and the dedicated guide \emph{The Complete Guide to Building Production-Grade
ML Projects That Get You Hired}. Both make the same point from different angles: deployment alone does not make a
project production-grade, and the gap between a tutorial project and a production one is what recruiters and
founders are actually screening for.

\section{Why Projects Get Rejected}

\subsection{What Recruiters Actually See}

\begin{example}
\textbf{Your resume:} ``Built ML model for house price prediction achieving 95\% accuracy using Random Forest.''

\textbf{Their reaction:} Skip. Next candidate.
\end{example}

Why? Because hundreds of other resumes say the exact same thing. That sentence signals: you followed a Kaggle
tutorial, you don't understand production systems, you'll need months to become useful, and you think accuracy is
what matters.

\subsection{What Recruiters Want to See}

Projects that answer four questions:

\begin{itemize}
  \item \textbf{Can you structure ambiguous problems?} Did you define success before building anything? Do you
  understand the business context or just the ML part? Can you explain why this problem needs ML and not a SQL
  query?
  \item \textbf{Can you build systems that don't break?} Did you think about what happens when data changes? How
  do you handle edge cases and failures? Is there a plan for when predictions are wrong?
  \item \textbf{Can you make technical decisions?} Why this model and not the others you tried? What's the
  trade-off between complexity and maintainability? Did you consider deployment constraints before choosing
  tools?
  \item \textbf{Can you communicate impact?} What business metric improved, and by how much? What does it cost
  to run this system? How do you know it's working in production?
\end{itemize}

\subsection{Tutorial Projects vs.\ Production Projects}

\begin{longtable}{p{0.47\textwidth} p{0.47\textwidth}}
\toprule
\textbf{Tutorial project} & \textbf{Production project} \\
\midrule
\endhead
Clean data provided & Messy data you have to deal with \\
Single metric to optimize & Multiple competing objectives \\
Train model, get score, done & Ongoing iteration and improvement \\
No consideration of deployment & Deployment is half the work \\
Works on your laptop & Works reliably at scale \\
\bottomrule
\end{longtable}

\begin{keytakeaway}
The projects that get you hired are the ones that close this gap --- not the ones with the fanciest model.
\end{keytakeaway}

\section{What a Production-Grade Project Must Demonstrate}

\begin{longtable}{p{0.24\textwidth} p{0.68\textwidth}}
\toprule
\textbf{Component} & \textbf{What you should demonstrate} \\
\midrule
\endhead
Problem definition & You understand the actual user or business problem \\
AI utilization & AI tools or agents accelerate appropriate parts of development \\
Architecture & You personally understand and can justify the important design decisions \\
Deployment & The system can be run and accessed outside your development environment \\
Monitoring & You can observe failures, performance, and operational behaviour \\
Testing & You can evaluate whether the system works correctly \\
Failure handling & You understand how the system breaks and what happens when it does \\
Debugging & You can reproduce, diagnose, and fix failures \\
Documentation & You explain decisions, limitations, failures, and fixes \\
Cost awareness & You understand the computational or inference cost where relevant \\
\bottomrule
\end{longtable}

\subsection{How to Present the Project}

Don't just publish the finished app --- document the engineering process: state the problem and why it matters;
explain the architecture and the decisions you made; identify where AI was used and where you exercised
independent judgment; describe a failure, unexpected result, or limitation; explain how you diagnosed and fixed
it; show evidence the improved version works.

\begin{example}
\bad \textbf{Weak:} ``Built an AI-powered application using Python and machine learning.''

\good \textbf{Stronger:} ``Built and deployed an AI-powered document-processing system, evaluated its failure
cases, optimized its processing pipeline, and documented reliability and cost trade-offs.''

An even stronger claim adds \emph{measured} results (reduced inference cost, improved latency) --- but only when
those results were actually measured and can be demonstrated, never invented.
\end{example}

\section{Choosing Problems That Prove You Can Think}

\subsection{The Problem Selection Framework}

\textbf{Rule 1 --- real problems only.} Stop with Titanic, Iris, MNIST, and clean Kaggle datasets. A problem is
``real'' when: the data has quality issues you must handle; there is clear business value you can articulate;
there is a deployment scenario that constrains your choices; there are ongoing maintenance requirements; multiple
stakeholders have different needs.

\begin{example}
\good \textbf{Good problem:} ``Predict restaurant demand to reduce food waste by 20\%, accounting for weather,
events, and seasonal patterns, deployed as daily batch predictions reviewed by managers.''

\bad \textbf{Bad problem:} ``Predict house prices on the Boston housing dataset.''
\end{example}

\textbf{Rule 2 --- production complexity required.} The project should force you to deal with missing data and
outliers, class imbalance or distributional shift, feature engineering from raw data, model-selection trade-offs,
serving-latency constraints, and monitoring/debugging needs. If the project doesn't force hard decisions, it
won't impress anyone.

\textbf{Rule 3 --- clear business context.} Answer before starting: What's the baseline (current solution --- a
heuristic, manual process, competitor product)? If you can't beat a three-line SQL query, don't use ML. What's the
success metric (translate ML metrics --- accuracy, F1 --- into business metrics: revenue, cost, time, retention)?
What's the deployment scenario (batch, real-time API, streaming)? What's the failure mode and its cost (a missed
diagnosis is catastrophic; a missed email is merely annoying)? What's the data refresh strategy (daily, weekly, or
only on performance degradation)?

\subsection{Project Categories That Work}

\begin{description}
  \item[Local business problems] A real business near you --- demand forecasting for restaurants, inventory
  optimization for retail, churn prediction for gyms, no-show prediction for clinics. Real consequences, real
  users to validate with, clear ROI.
  \item[Personal pain points] A problem you actually face --- a job-application tracker with outcome prediction,
  a reading-habit recommender, a study-schedule optimizer, expense categorization. You understand it deeply and
  can iterate on real usage.
  \item[Recreate real systems] Rebuild a feature from a major product --- YouTube thumbnail click prediction,
  Spotify playlist generation, Netflix recommendations, Twitter trend detection, Amazon review-helpfulness
  ranking. Well-understood problem space you can benchmark against.
  \item[Domain-specific applications] Pick an industry you want to work in --- healthcare (readmission risk,
  image analysis, drug-interaction detection), finance (credit default with fairness constraints, fraud
  detection, sentiment analysis), e-commerce (cold-start recommendations, dynamic pricing, search ranking). Shows
  domain interest and gets you interviews in that domain.
\end{description}

\subsection{Selection Criteria Checklist}

Before committing to a project, verify: \textbf{data accessibility} (can you legally and ethically get this
data?); \textbf{scope appropriateness} (finishable in 2--4 weeks --- not a day, not six months; one well-done
project beats five half-finished ones); \textbf{demonstrability} (explainable to a recruiter in two minutes --- if
it takes twenty, it's too complex); \textbf{learning value} (does it teach you something new?); \textbf{portfolio
fit} (does it complement your other projects --- aim for one end-to-end pipeline, one deployment-focused build,
one domain-specific build?).

\section{The Production Mindset vs.\ the Student Mindset}

\begin{keytakeaway}
\textbf{Student mindset:} ``How do I get the highest accuracy on this test set?'' \\
\textbf{Production mindset:} ``How do I build a system that reliably delivers business value over time?''
\end{keytakeaway}

\begin{longtable}{p{0.18\textwidth} p{0.38\textwidth} p{0.38\textwidth}}
\toprule
\textbf{Decision} & \textbf{Student mindset} & \textbf{Production mindset} \\
\midrule
\endhead
Model selection & Pick the most complex model that gives the best test score & Pick the simplest model that meets
business requirements --- easier to debug, explain, retrain, and more robust to shift \\
Feature engineering & Create 100 features, let the model figure it out & Create interpretable features you can
explain and monitor \\
Evaluation & Single train--test split, report F1 & Time-based splits, multiple metrics, business-impact analysis
\\
Deployment & Model works on laptop, job done & Containerized, monitored, versioned, with a rollback plan \\
Documentation & Comments in a notebook, if anything & README, API docs, deployment guide, monitoring playbook \\
\bottomrule
\end{longtable}

\subsection{The Questions Production Systems Must Answer}

For every component you build: \textbf{What happens when this breaks?} (Is there a fallback? How is the failure
detected? What's the recovery process? Who gets alerted?) \textbf{How do you know it's working?} (What metrics
indicate health? What's the monitoring strategy? What triggers investigation?) \textbf{How does this change over
time?} (Will data distribution shift? Will requirements or scale grow? Is the system adaptable?) \textbf{What's
the maintenance burden?} (How often does it need updating? What skills does maintaining it require? Is this
sustainable?)

\section{Phase 1 --- Problem Definition \& Business Context}

\subsection{Step 1: Write a Project Brief}

Before writing any code, document: \textbf{the business problem} (not ``predict X'' but ``reduce waste / increase
revenue / save time by doing X''); \textbf{the current baseline} (manual process --- quantify the time spent;
simple heuristic --- quantify its performance; competitor product --- understand its limitations; or, if
unsolved, quantify the cost of that); \textbf{the success definition}, translating ML metrics into business
metrics from day one (``5\% accuracy increase'' is meaningless; ``10\% cost reduction'' is a concrete dollar
figure).

\textbf{Constraints \& requirements} to pin down up front: time constraints (predictions needed daily at 6\,a.m.?
real-time under 100\,ms? --- this determines the whole architecture); cost constraints (free-tier tools only?
budget ceiling?); accuracy requirements (99\% for healthcare/finance vs.\ 80\% acceptable for recommendations);
explainability requirements (regulated industries must explain every prediction); data constraints (privacy
regulations such as GDPR/HIPAA, data availability, labeling resources).

\subsection{Step 2: Define a Success-Metrics Hierarchy}

\begin{description}
  \item[Primary business metric] The one number stakeholders care about most --- revenue impact, cost savings,
  time saved, user satisfaction.
  \item[Secondary ML metrics] The metrics you optimize during development --- precision, recall, RMSE, MAE ---
  chosen based on business context.
  \item[Guardrail metrics] The metrics you cannot sacrifice --- fairness, privacy, latency, interpretability.
\end{description}

\begin{example}
\textbf{Example hierarchy --- customer churn prediction for a subscription service:}
Primary: monthly revenue retention (\$). Secondary: recall on churners (catch 80\%+), precision (limit false
alarms). Guardrails: no demographic bias; predictions explainable to the customer-success team.

This matters because when accuracy and fairness conflict, you now have a clear priority; when speed and precision
conflict, you know what to optimize.
\end{example}

\subsection{Step 3: Identify Stakeholders \& Their Needs}

\begin{itemize}
  \item \textbf{End users} --- who receives predictions, in what format (API, dashboard, email), at what
  frequency, and do they need explainability?
  \item \textbf{Decision makers} --- who approves the project; they care about ROI justification, risk
  assessment, resource requirements, and timeline.
  \item \textbf{Maintainers} --- who keeps it running; they need documentation, monitoring dashboards, alert
  systems, and deployment runbooks.
  \item \textbf{Data providers} --- who owns the data; they care about privacy requirements, access patterns,
  update frequency, and quality guarantees.
\end{itemize}

\subsection{Step 4: Risk Assessment}

\begin{longtable}{p{0.22\textwidth} p{0.72\textwidth}}
\toprule
\textbf{Risk category} & \textbf{Examples} \\
\midrule
\endhead
Data risks & Source becomes unavailable (vendor shutdown, API deprecated); quality degrades over time (logging
errors, schema changes); distribution shifts (market or user-behaviour changes); privacy violations (accidentally
logging PII) \\
Model risks & Performance degrades silently (undetected drift); adversarial inputs (gaming the system); bias
amplification; edge cases not covered in training \\
System risks & Downtime during critical periods; scaling failures under load; dependency breakage; security
vulnerabilities \\
Business risks & Incorrect predictions drive bad decisions; over-reliance on automation; regulatory compliance
issues; reputational damage from failures \\
\bottomrule
\end{longtable}

For each risk, define its likelihood (low/medium/high), impact (minimal/moderate/severe), mitigation strategy, and
monitoring approach. Recruiters routinely ask ``what could go wrong with this system?'' --- have the answer ready.

\section{Phase 2 --- Data Strategy \& Pipeline Design}

\begin{keytakeaway}
Raw data is sacred --- never modify originals. Pipeline stages: Raw $\rightarrow$ Validated $\rightarrow$
Processed $\rightarrow$ Features $\rightarrow$ Model-ready. Each stage is versioned, reproducible, and documented.
\end{keytakeaway}

\subsection{Step 1: Data Acquisition Strategy}

\textbf{Primary sources} --- internal systems (databases, logs, APIs), external APIs (weather, market data, public
datasets), web scraping (respect robots.txt, rate limits, terms of service), or manual collection (surveys,
annotations). For each source, document its access method, update frequency, cost, reliability, and legal
constraints. \textbf{Supplementary sources} enrich the primary data --- weather for demand forecasting, holidays
for sales prediction, economic indicators for financial models, demographics for personalization.

\begin{caution}
\textbf{Data-collection ethics.} Before scraping or collecting data, verify: robots.txt allows it; the terms of
service permit it; privacy laws don't prohibit it (GDPR, CCPA); rate limiting won't get you blocked; attribution
and commercial-use restrictions are respected. If you can't answer these, don't use that data --- recruiters will
ask about data ethics, and many candidates fail exactly here.
\end{caution}

\subsection{Step 2: Data Quality Framework}

Build validation checks across five dimensions, not just ``check for nulls'': \textbf{completeness} (what
percentage of missing values is acceptable; which fields are critical); \textbf{validity} (valid ranges, formats,
consistency rules); \textbf{accuracy} (how you validate correctness and how often you audit); \textbf{consistency}
(do related fields and units agree); \textbf{timeliness} (how fresh must the data be, and how do you detect
staleness).

\subsection{Step 3: Data Versioning Strategy}

Version data for \textbf{reproducibility} (must reproduce any experiment exactly --- ``trained on March 15'' means
nothing without a data version), \textbf{debugging} (know which data a failing production model was trained on),
\textbf{auditability} (prove which data was used, for regulated industries), and \textbf{experimentation}
(compare performance across data versions). Tools: DVC (industry standard), Git LFS (smaller datasets), or custom
hashing for minimal setups. Version raw snapshots, processed outputs, feature outputs, train/validation/test
splits, and the processing scripts themselves --- with metadata for collection timestamp, source, pipeline
version, quality metrics, and known issues.

\subsection{Step 4: Feature Engineering Strategy}

\begin{example}
\bad \textbf{Students:} create random features, hope something works.

\good \textbf{Professionals:} create features based on explicit hypotheses about the problem.
\end{example}

Process: (1) understand the problem domain --- what do domain experts consider important (seasonality and
promotions for sales; velocity and historical behaviour for fraud; user preferences and context for
recommendations); (2) form a hypothesis (``day of week matters because weekend behaviour differs''); (3) create
features that test the hypothesis; (4) validate feature importance --- did it help, by how much, was the
hypothesis correct; (5) document the reasoning.

Feature categories to consider: \textbf{temporal} (trends, seasonality, recency, frequency), \textbf{aggregations}
(statistics over time windows, counts, ratios), \textbf{interactions} (non-linear combinations, e.g.\ price
$\times$ quantity $=$ revenue), \textbf{encodings} (target/frequency encoding for categoricals, TF-IDF/embeddings
for text, pretrained-model features for images). The key rule: every feature should be explainable and
monitorable --- if you can't explain why it matters, don't include it; if you can't monitor it for drift, it's a
risk.

\subsection{Step 5: Data Pipeline Architecture}

Design principles: \textbf{modularity} (preprocessing $\rightarrow$ validation $\rightarrow$ feature engineering
$\rightarrow$ split, each independently testable); \textbf{reproducibility} (same input always produces the same
output --- lock random seeds, version dependencies, save processing parameters); \textbf{observability} (log
input/output shapes, statistics, timing, errors at each step); \textbf{error handling} (define behaviour for
missing source data, validation failures, processing errors --- alert-and-fallback, stop-and-alert, or
log-and-continue); \textbf{scalability} (batch for large datasets, streaming for real-time needs, caching for
expensive operations). The pipeline should be automated, scheduled, monitored, and documented.

